% !TeX root = ../../Book2.tex
% 纲要标题：### 17.4 具体代数
% 本轮补充的纲要细目：
% * 以四元数左右乘构造范数相似变换，证明乘子及对角标量核；不由核的计算断言满射
% 纲要位置：计划纲要.md，第 1325 行。

\section{具体代数}
\label{b2:sec:1704}

四元数代数是次数二的中心单代数模型。它的乘法只需两个参数就能写下，是否分裂却与一个范数方程的可解性有关。本节假设 \(\operatorname{char}k\ne2\)，并固定 \(a,b\in k^\times\)；这一假设用于基的验证、共轭及二次扩张的描述。

% 纲要要点（供目录核对）：
%
% * 四元数代数
% * 矩阵代数及非分裂例子
% * 与二次型及范数方程的联系
%
% ---
%

\subsection{四元数代数}
\label{b2:subsec:1704:01}

\begin{definition}\label{b2:def:quaternion-algebra}
设 \(k\) 为特征不等于 \(2\) 的域，\(a,b\in k^\times\)。将自由结合代数 \(k\langle i,j\rangle\) 对下列关系生成的双边理想取商：
\[
i^2=a,\qquad j^2=b,\qquad ji=-ij
\]
所得含幺 \(k\) 代数称为\term[siyuanshu-daishu]{四元数代数}[quaternion algebra]，记为 \((a,b)_k\)。
\end{definition}
\symidx{\((a,b)_k\)：参数为 \(a,b\) 的四元数代数}

关系使每个字都能先把 \(i\) 移到 \(j\) 左边，再将两个指数降到零或一，所以 \(1,i,j,ij\) 生成这个代数。它们还线性无关，这一点不能仅由关系写法直接断言。

\begin{proposition}\label{b2:prop:quaternion-basis-splitting}
设 \(k\) 为特征不等于 \(2\) 的域，\(a,b\in k^\times\)，\(Q=(a,b)_k\)，则 \(Q\) 以 \(1,i,j,ij\) 为 \(k\) 基，并且是次数二的中心单 \(k\) 代数。若域扩张 \(K/k\) 含有 \(\alpha\) 满足 \(\alpha^2=a\)，则有分裂同构
\[
Q_K=K\otimes_kQ\longrightarrow M_2(K),\qquad
i\longmapsto I=\begin{pmatrix}\alpha&0\\0&-\alpha\end{pmatrix},
\quad j\longmapsto J=\begin{pmatrix}0&b\\1&0\end{pmatrix}.
\]
\end{proposition}
\begin{proof}
这两个矩阵满足 \(I^2=aI_2\)、\(J^2=bI_2\)、\(JI=-IJ\)，所以给出由商代数到矩阵代数的同态。在对角位置，\(I_2,I\) 线性无关，因为 \(2\alpha\ne0\)；在非对角位置，\(J,IJ\) 线性无关，因为其两个坐标组成的行列式为 \(-2\alpha b\ne0\)。故四个矩阵为 \(M_2(K)\) 的基。原四个生成元若有 \(k\) 线性关系，映射后也有同一关系，因此系数全零。这证明 \(\dim_kQ=4\)，扩张后同态便是基之间的同构。

现在验证 \(Q\) 本身中心单。若 \(U\) 为非零双边理想，则 \(K\otimes_kU\) 是 \(Q_K\cong M_2(K)\) 的非零理想，因矩阵环单而等于整个 \(Q_K\)。维数比较得 \(\dim_kU=4\)，所以 \(U=Q\)。若 \(z\in Z(Q)\)，其像在 \(M_2(K)\) 中为标量。以 \(I_2,I,J,IJ\) 为基比较系数，得到 \(z\) 只能是原 \(1\) 的 \(k\) 倍数，故中心为 \(k\)。维数为四使次数为二。
\end{proof}

对 \(q=x_0+x_1i+x_2j+x_3ij\)，定义它的\term[siyuanshu-gonge]{四元数共轭}[quaternion conjugation]为
\[
\bar q=x_0-x_1i-x_2j-x_3ij.
\]
生成元的像 \(i\mapsto-i,j\mapsto-j\) 在反代数中仍满足原关系，所以该映射反转乘法；它在上述基上作用两次为恒等，因此
\begin{equation}\label{b2:eq:quaternion-conjugation}
\overline{qr}=\bar r\bar q,\qquad\overline{\bar q}=q.
\end{equation}
直接按乘法关系展开，交叉项成对消去，得到
\begin{equation}\label{b2:eq:quaternion-norm}
q\bar q=\bar q q
=x_0^2-a x_1^2-b x_2^2+ab x_3^2.
\end{equation}

\begin{proposition}\label{b2:prop:quaternion-norm}
设 \(k\) 为特征不等于 \(2\) 的域，\(a,b\in k^\times\)，\(Q=(a,b)_k\)。对 \(q=x_0+x_1i+x_2j+x_3ij\in Q\)，令 \(\bar q=x_0-x_1i-x_2j-x_3ij\)。由于
\[
q\bar q=\bar q q=x_0^2-a x_1^2-b x_2^2+ab x_3^2\in k,
\]
可将这个标量记为 \(N(q)\)，称为\term[siyuanshu-fanshu]{四元数范数}[quaternion norm]。对任意 \(q,r\in Q\)，有
\[
N(qr)=N(q)N(r),\qquad
q\in Q^\times\Longleftrightarrow N(q)\ne0.
\]
可逆时 \(q^{-1}=N(q)^{-1}\bar q\)。在\cref{b2:prop:quaternion-basis-splitting}的分裂矩阵模型中，\(N(q)\) 就是矩阵行列式。
\end{proposition}
\begin{proof}
因 \(N(r)\in k\) 在中心，
\[
N(qr)=qr\bar r\bar q=qN(r)\bar q=N(q)N(r).
\]
若范数非零，\eqref{b2:eq:quaternion-norm}给出所述双侧逆。若 \(q\) 可逆，则乘法性给出 \(N(q)N(q^{-1})=N(1)=1\)，所以范数非零。

在分裂模型中，\(q\) 的矩阵为
\[
\begin{pmatrix}
x_0+\alpha x_1&b(x_2+\alpha x_3)\\
x_2-\alpha x_3&x_0-\alpha x_1
\end{pmatrix}.
\]
其行列式为 \(x_0^2-a x_1^2-bx_2^2+abx_3^2\)，正好是 \(N(q)\)。
\end{proof}
\symidx{\(N(q)=q\bar q\)：四元数范数}

\begin{proposition}\label{b2:prop:quaternion-norm-similitudes}
设 \(k\) 为特征不为二的域，\(Q=(a,b)_k\)、\(a,b\in k^\times\)，\(N(x)=x\bar x\) 为四元数范数。对 \(u,v\in Q^\times\)，定义 \(k\) 线性映射
\[
L_{u,v}:Q\longrightarrow Q,\qquad x\longmapsto uxv^{-1}.
\]
它是范数二次型的相似变换，乘子为 \(N(u)/N(v)\)。映射 \((u,v)\mapsto L_{u,v}\) 是从 \(Q^\times\times Q^\times\) 到该相似变换群的群同态，其核为 \(\{(c,c)\mid c\in k^\times\}\)。特别地，内自同构 \(x\mapsto uxu^{-1}\) 保持 \(N\)。
\end{proposition}
\begin{proof}
范数的对角系数为 \(1,-a,-b,ab\)，均非零，所以二次型非退化。\cref{b2:prop:quaternion-norm}的乘法性给出
\[
N(uxv^{-1})=\frac{N(u)}{N(v)}N(x).
\]
由特征不为二，极化同样使相应双线性型乘以这个标量。又有
\[
L_{u,v}L_{u',v'}(x)=uu'xv'^{-1}v^{-1}
=L_{uu',vv'}(x),
\]
故是群同态。若 \(L_{u,v}\) 为恒等，在 \(x=1\) 处即得 \(u=v\)；再对全部 \(x\) 使用等式，得到 \(u\) 属于 \(Q\) 的中心，也就是 \(k^\times\)。反过来，对角标量显然作用为恒等。这证明核的描述，而取 \(u=v\) 就得到等距结论。
\end{proof}

这给出一族由代数乘法产生的几何变换。仅从核的计算不能断言它们穷尽全部相似变换；这里只使用上述明确构造，也不预先调用正交群的进一步分类。

\subsection{矩阵代数与非分裂例子}
\label{b2:subsec:1704:02}

\begin{proposition}\label{b2:prop:quaternion-division-or-split}
设 \(k\) 为特征不等于 \(2\) 的域，\(a,b\in k^\times\)，\(Q=(a,b)_k\)，\(N:Q\rightarrow k\) 为其四元数范数，则 \(Q\) 或者是除代数，或者同构于 \(M_2(k)\)。它分裂，当且仅当二次型 \(N\) 有非平凡零点，即存在 \(q\in Q\setminus\{0\}\) 使 \(N(q)=0\)。
\end{proposition}
\begin{proof}
由\cref{b2:prop:csa-division-description}，\(Q\cong M_r(D)\)，其中 \(D\) 是除代数，且 \(4=r^2\dim_kD\)。若 \(r=1\)，\(Q\) 是除代数；否则只能 \(r=2,\dim_kD=1\)，即 \(D=k\)，所以 \(Q\cong M_2(k)\)。

由\cref{b2:prop:quaternion-norm}，每个非零元素可逆当且仅当每个非零元素的范数都非零。因此有非平凡零点恰好等价于不是除代数，按前一分类又等价于分裂。
\end{proof}

若 \(a=\alpha^2\in k^{\times2}\)，\cref{b2:prop:quaternion-basis-splitting}的两个矩阵已经定义在 \(k\) 中，因此 \((a,b)_k\) 直接分裂。同样，交换生成元 \(i,j\) 给出 \((a,b)_k\cong(b,a)_k\)，所以 \(b\) 为平方时也分裂。

\ProseParagraphBreak
实四元数代数
\[
\mathbb H=(-1,-1)_{\mathbb R}
\]
则不分裂，因为其范数为 \(x_0^2+x_1^2+x_2^2+x_3^2\)，对任意非零实向量严格为正。因而每个非零四元数都可逆；矩阵代数 \(M_2(\mathbb R)\) 却有非零奇异矩阵，二者不可能同构。扩张到 \(\mathbb C\) 后，参数 \(-1\) 成为平方，所以 \(\mathbb H\otimes_{\mathbb R}\mathbb C\cong M_2(\mathbb C)\)。

\ProseParagraphBreak
更一般地，对非零实数 \(a,b\)，只要其中一个为正，就因该参数为平方而分裂；若二者都为负，分别把 \(i,j\) 除以 \(\sqrt{-a},\sqrt{-b}\)，得到平方均为 \(-1\) 的生成元，故代数同构于 \(\mathbb H\)。因此实四元数代数恰有分裂矩阵代数与 \(\mathbb H\) 两种类型。这里分类的是四维四元数代数，没有在此断言所有实中心单代数的进一步分类。

\subsection{二次型与范数方程}
\label{b2:subsec:1704:03}

\begin{theorem}[四元数的范数分裂判据]\label{b2:thm:quaternion-splitting}
设 \(k\) 为特征不等于 \(2\) 的域，\(a,b\in k^\times\)。若 \(a\) 不是 \(k\) 中的平方，令 \(L=k(\alpha)\)，其中 \(\alpha^2=a\)，则
\[
(a,b)_k\text{ 分裂}\quad\Longleftrightarrow\quad
b\in N_{L/k}(L^\times)
\quad\Longleftrightarrow\quad b=x^2-a y^2\text{ 对某些 }x,y\in k.
\]
若 \(a\) 为非零平方，代数总分裂，最后一个方程也总有解。
\end{theorem}
\begin{proof}
\(a\) 非平方时，\(L=k[i]\subseteq Q\) 是二次域，且 \(Q=L\oplus Lj\)。记 \(L\) 的非平凡自同构为 \(u\mapsto\bar u\)，则 \(ju=\bar u j\)。由前述范数公式，对 \(u,v\in L\) 有
\[
N(u+vj)=N_{L/k}(u)-bN_{L/k}(v).
\]
这里 \(N_{L/k}(x+y\alpha)=(x+y\alpha)(x-y\alpha)=x^2-a y^2\)，可由二次扩张的乘法矩阵或两个共轭直接计算。

若 \(Q\) 分裂，\cref{b2:prop:quaternion-division-or-split}给出非零 \(u+vj\) 的范数为零。不能有 \(v=0\)，否则域范数零使 \(u=0\)，矛盾。因此 \(v\ne0\)，得到
\(b=N_{L/k}(u)/N_{L/k}(v)=N_{L/k}(u/v)\)。反过来，若 \(b=N_{L/k}(z)\)，则 \(z+j\ne0\) 且其范数为零，所以代数分裂。域范数的显式公式又给出最后一个等价。

若 \(a=\alpha^2\) 且 \(\alpha\in k^\times\)，分裂性已证明，而
\[
x=\frac{b+1}{2},\qquad y=\frac{b-1}{2\alpha}
\]
满足 \(x^2-a y^2=b\)，故方程仍可解。
\end{proof}

已知范数方程的一个解，还可以直接构造分裂同构。若 \(a\) 非平方，取 \(z=x+y\alpha\in L\) 满足 \(N_{L/k}(z)=b\)。在二维 \(k\) 空间 \(L\) 上，让 \(i\) 作用为乘 \(\alpha\)，让 \(j\) 作用为 \(v\mapsto z\bar v\)。两者满足平方关系，且
\(z\overline{\alpha v}=-\alpha z\bar v\)，所以反交换，给出 \(Q\to\operatorname{End}_k(L)\)。它保持单位元，由单性为单射，两端维数同为四，故同构。在基 \(1,\alpha\) 下，具体矩阵为
\[
i\longmapsto\begin{pmatrix}0&a\\1&0\end{pmatrix},\qquad
j\longmapsto\begin{pmatrix}x&-ay\\y&-x\end{pmatrix}.
\]
后者的平方为 \((x^2-a y^2)I_2=bI_2\)，可作为公式的直接核验。

\ProseParagraphBreak
范数判据也等价于齐次方程
\[
X^2-aY^2-bZ^2=0
\]
有一个非零解。若 \(Z\ne0\)，除以 \(Z^2\) 就得到 \(b=(X/Z)^2-a(Y/Z)^2\)；若 \(Z=0\)，非零解必有 \(Y\ne0\)，于是 \(a=(X/Y)^2\)，同样推出分裂。反过来，范数方程解 \(b=x^2-a y^2\) 给出点 \((x,y,1)\)，平方参数时也可直接取 \((\alpha,1,0)\)。这样，四元数代数是否为矩阵代数，已经化为一个具体二次方程的可解性问题。
